% chap04_functions.tex - Defining and using functions
% ============================================================================

\chapter{Functions}
\label{chap:functions}

You've been using functions like \py{print()}, \py{len()}, and \py{type()} throughout this book. Now it's time to write your own. Functions let you package code into reusable pieces, making your programs shorter, clearer, and easier to maintain.

% ----------------------------------------------------------------------------
\section{Why Functions?}
\label{sec:why-functions}

Imagine you need to calculate molarity in multiple places in your code:

\begin{lstlisting}
# First calculation
mass1 = 5.85
molar_mass1 = 58.44
volume1 = 0.5
molarity1 = mass1 / molar_mass1 / volume1

# Second calculation
mass2 = 11.7
molar_mass2 = 58.44
volume2 = 1.0
molarity2 = mass2 / molar_mass2 / volume2

# Third calculation...
\end{lstlisting}

This is repetitive. If you later realize you made an error in the formula, you'd need to fix it in every location. Functions solve this by letting you define the calculation once and use it many times.

% ----------------------------------------------------------------------------
\section{Defining Functions}
\label{sec:defining-functions}

Create a function with the \py{def} keyword:

\begin{lstlisting}
def greet():
    print("Hello, scientist!")
\end{lstlisting}

The structure is:
\begin{itemize}
    \item \py{def} keyword
    \item Function name (follows same rules as variable names)
    \item Parentheses (may contain parameters)
    \item Colon
    \item Indented block of code (the function body)
\end{itemize}

To use (``call'') the function:

\begin{lstlisting}
greet()  # prints "Hello, scientist!"
greet()  # prints it again
\end{lstlisting}

\subsection{Parameters and Arguments}

Functions become powerful when they accept inputs called \emph{parameters}:

\begin{lstlisting}
def greet(name):
    print(f"Hello, {name}!")

greet("Alice")   # prints "Hello, Alice!"
greet("Bob")     # prints "Hello, Bob!"
\end{lstlisting}

In the definition, \py{name} is a \emph{parameter}---a placeholder. When you call the function, you provide an \emph{argument}---the actual value. (Many people use these terms interchangeably, but technically parameters are in definitions, arguments are in calls.)

Multiple parameters are separated by commas:

\begin{lstlisting}
def calculate_molarity(mass, molar_mass, volume):
    molarity = mass / molar_mass / volume
    print(f"Molarity: {molarity:.4f} M")

calculate_molarity(5.85, 58.44, 0.5)
calculate_molarity(11.7, 58.44, 1.0)
\end{lstlisting}

\subsection{Return Values}

Printing is nice, but usually you want to get a value back from a function. Use \py{return}:

\begin{lstlisting}
def calculate_molarity(mass, molar_mass, volume):
    molarity = mass / molar_mass / volume
    return molarity

# Now you can use the result
conc = calculate_molarity(5.85, 58.44, 0.5)
print(f"Concentration: {conc:.4f} M")

# Or use it directly in expressions
if calculate_molarity(11.7, 58.44, 1.0) > 0.1:
    print("Concentration is above threshold")
\end{lstlisting}

\begin{note}
A function without a \py{return} statement returns \py{None}. If you try to use the result, you'll get \py{None}---a common source of confusion for beginners.
\end{note}

\subsection{Multiple Return Values}

Python functions can return multiple values using tuples:

\begin{lstlisting}
def analyze_data(values):
    minimum = min(values)
    maximum = max(values)
    average = sum(values) / len(values)
    return minimum, maximum, average

data = [1.2, 3.4, 2.1, 5.6, 4.3]
low, high, avg = analyze_data(data)
print(f"Range: {low} to {high}, Average: {avg:.2f}")
\end{lstlisting}

\begin{labexample}
A function to calculate reaction quotient Q:

\begin{lstlisting}
def reaction_quotient(products, reactants):
    """
    Calculate Q for a reaction.
    products: list of [concentration, coefficient] pairs
    reactants: list of [concentration, coefficient] pairs
    """
    Q_num = 1.0
    for conc, coeff in products:
        Q_num *= conc ** coeff

    Q_den = 1.0
    for conc, coeff in reactants:
        Q_den *= conc ** coeff

    return Q_num / Q_den

# For: 2NO2 <-> N2O4
# [N2O4] = 0.5 M, [NO2] = 0.1 M
Q = reaction_quotient(
    products=[[0.5, 1]],      # N2O4 with coefficient 1
    reactants=[[0.1, 2]]      # NO2 with coefficient 2
)
print(f"Q = {Q}")
\end{lstlisting}
\end{labexample}

% ----------------------------------------------------------------------------
\section{Default Parameters}
\label{sec:default-params}

You can give parameters default values:

\begin{lstlisting}
def calculate_molarity(mass, molar_mass, volume=1.0):
    return mass / molar_mass / volume

# volume defaults to 1.0 if not specified
conc1 = calculate_molarity(5.85, 58.44)        # volume = 1.0
conc2 = calculate_molarity(5.85, 58.44, 0.5)   # volume = 0.5
\end{lstlisting}

\begin{warning}
Parameters with defaults must come after parameters without defaults:
\begin{lstlisting}
def func(a, b=1, c):    # ERROR! c has no default but comes after b
    pass

def func(a, c, b=1):    # OK - default parameter is last
    pass
\end{lstlisting}
\end{warning}

\subsection{The Mutable Default Argument Trap}
\label{subsec:mutable-defaults}

One of Python's most notorious gotchas involves mutable default arguments. Consider this function:

\begin{lstlisting}
def add_reading(reading, readings=[]):
    readings.append(reading)
    return readings

# Seems fine at first...
print(add_reading(1.5))   # [1.5]
print(add_reading(2.3))   # [1.5, 2.3]  -- wait, what?
print(add_reading(3.1))   # [1.5, 2.3, 3.1]  -- it remembers!
\end{lstlisting}

The list accumulates values across calls! This happens because the default value \py{[]} is created \emph{once} when the function is defined, not each time it's called. Every call shares the same list object.

Recall from Chapter~\ref{chap:basics} (Section~\ref{subsec:aliasing}) that variables are labels pointing to objects. The default \py{[]} creates one list object, and every call that uses the default gets that same object.

\begin{tip}
The fix is to use \py{None} as the default and create the list inside the function:
\begin{lstlisting}
def add_reading(reading, readings=None):
    if readings is None:
        readings = []    # create a NEW list each time
    readings.append(reading)
    return readings

print(add_reading(1.5))   # [1.5]
print(add_reading(2.3))   # [2.3] -- fresh list each time
\end{lstlisting}
This pattern applies to any mutable default: lists, dictionaries, sets, or custom objects.
\end{tip}

% ----------------------------------------------------------------------------
\section{Keyword Arguments}
\label{sec:keyword-args}

When calling functions, you can specify arguments by name:

\begin{lstlisting}
def describe_element(name, symbol, atomic_number):
    print(f"{name} ({symbol}): Z = {atomic_number}")

# Positional arguments (order matters)
describe_element("Carbon", "C", 6)

# Keyword arguments (order doesn't matter)
describe_element(atomic_number=6, name="Carbon", symbol="C")

# Mixed (positional first, then keyword)
describe_element("Carbon", atomic_number=6, symbol="C")
\end{lstlisting}

Keyword arguments make function calls more readable, especially with many parameters:

\begin{lstlisting}
# What do these numbers mean?
result = analyze_spectrum(data, 400, 700, 0.1, True, False)

# Much clearer:
result = analyze_spectrum(
    data,
    wavelength_min=400,
    wavelength_max=700,
    resolution=0.1,
    normalize=True,
    subtract_baseline=False
)
\end{lstlisting}

% ----------------------------------------------------------------------------
\section{Flexible Arguments: *args and **kwargs}
\label{sec:args-kwargs}

When reading library documentation, you'll often see function signatures like:

\begin{lstlisting}
def plot(x, y, *args, **kwargs):
    ...
\end{lstlisting}

What do the asterisks mean?

\subsection{*args: Variable Positional Arguments}

A parameter prefixed with \py{*} collects any extra positional arguments into a tuple:

\begin{lstlisting}
def calculate_mean(*values):
    """Calculate mean of any number of values."""
    return sum(values) / len(values)

print(calculate_mean(1, 2, 3))           # 2.0
print(calculate_mean(10, 20, 30, 40))    # 25.0
\end{lstlisting}

The name \py{args} is a convention---you could use any name, but \py{*args} is standard.

\subsection{**kwargs: Variable Keyword Arguments}

A parameter prefixed with \py{**} collects extra keyword arguments into a dictionary:

\begin{lstlisting}
def create_sample(name, **properties):
    """Create a sample with arbitrary properties."""
    print(f"Sample: {name}")
    for key, value in properties.items():
        print(f"  {key}: {value}")

create_sample("S1", concentration=0.5, temperature=25, pH=7.2)
# Output:
# Sample: S1
#   concentration: 0.5
#   temperature: 25
#   pH: 7.2
\end{lstlisting}

\subsection{Why This Matters}

You'll rarely need to define functions with \py{*args} and \py{**kwargs} yourself, but you'll \emph{call} many library functions that use them. Understanding this syntax helps you:

\begin{itemize}
    \item Read documentation for libraries like Matplotlib or Pandas
    \item Pass optional parameters through wrapper functions
    \item Understand error messages about unexpected arguments
\end{itemize}

\begin{lstlisting}
# Matplotlib's plot() accepts many optional keyword arguments
import matplotlib.pyplot as plt
plt.plot(x, y, color='red', linewidth=2, marker='o')
# These extra arguments are collected by **kwargs
\end{lstlisting}

% ----------------------------------------------------------------------------
\section{Docstrings}
\label{sec:docstrings}

A \emph{docstring} is a string at the beginning of a function that documents what it does:

\begin{lstlisting}
def calculate_molarity(mass, molar_mass, volume):
    """
    Calculate the molarity of a solution.

    Parameters:
        mass: Mass of solute in grams
        molar_mass: Molar mass of solute in g/mol
        volume: Volume of solution in liters

    Returns:
        Molarity in mol/L
    """
    return mass / molar_mass / volume
\end{lstlisting}

Docstrings help others (and future you) understand your code. Python's \py{help()} function displays them:

\begin{lstlisting}
>>> help(calculate_molarity)
\end{lstlisting}

\begin{tip}
Write docstrings for any function that isn't immediately obvious. It takes a minute now and saves hours of confusion later.
\end{tip}

% ----------------------------------------------------------------------------
\section{Scope: Where Variables Live}
\label{sec:scope}

Variables created inside a function exist only inside that function. This is called \emph{scope}:

\begin{lstlisting}
def my_function():
    x = 10  # x exists only inside this function
    print(x)

my_function()  # prints 10
print(x)       # ERROR! x is not defined here
\end{lstlisting}

Variables defined outside functions are ``global'' and can be \emph{read} from anywhere without special syntax:

\begin{lstlisting}
multiplier = 2  # global variable

def double(x):
    return x * multiplier  # reading global variable works fine

print(double(5))  # prints 10
\end{lstlisting}

However, \emph{modifying} a global variable from inside a function requires the \py{global} keyword:

\begin{lstlisting}
counter = 0

def increment_broken():
    counter = counter + 1  # ERROR! Python thinks counter is local

def increment_works():
    global counter         # tell Python we mean the global variable
    counter = counter + 1

increment_works()
print(counter)  # prints 1
\end{lstlisting}

Without \py{global}, Python assumes you're creating a new local variable. The error occurs because you're trying to read \py{counter} (on the right side of \py{=}) before assigning it locally.

\begin{warning}
Modifying global variables is generally discouraged---it makes code harder to understand and debug. Instead, pass values as parameters and return results:
\begin{lstlisting}
# Bad - modifies global state
counter = 0
def increment():
    global counter
    counter += 1

# Better - explicit input and output
def increment(counter):
    return counter + 1

counter = increment(counter)
\end{lstlisting}
Functions that don't modify global state are easier to test, reuse, and reason about.
\end{warning}

% ----------------------------------------------------------------------------
\section{Functions as Values}
\label{sec:functions-as-values}

In Python, functions are values---you can store them in variables, put them in lists, and pass them to other functions:

\begin{lstlisting}
def square(x):
    return x ** 2

def cube(x):
    return x ** 3

# Store function in a variable
operation = square
print(operation(4))  # prints 16

operation = cube
print(operation(4))  # prints 64

# Put functions in a list
operations = [square, cube]
for op in operations:
    print(op(3))  # prints 9, then 27
\end{lstlisting}

This is useful for applying different operations to data:

\begin{lstlisting}
def apply_to_all(data, function):
    """Apply a function to every element in a list."""
    return [function(x) for x in data]

numbers = [1, 2, 3, 4, 5]
squared = apply_to_all(numbers, square)  # [1, 4, 9, 16, 25]
\end{lstlisting}

% ----------------------------------------------------------------------------
\section{Lambda Functions}
\label{sec:lambda}

For simple, one-line functions, Python offers \py{lambda}---a way to create small anonymous (unnamed) functions:

\begin{lstlisting}
# Regular function
def square(x):
    return x ** 2

# Equivalent lambda
square = lambda x: x ** 2
\end{lstlisting}

Lambdas are useful when you need a quick function for a single use:

\begin{lstlisting}
numbers = [1, 2, 3, 4, 5]

# Using lambda with map
squared = list(map(lambda x: x ** 2, numbers))

# Using lambda for sorting
elements = [("C", 6), ("H", 1), ("O", 8), ("N", 7)]
sorted_by_number = sorted(elements, key=lambda x: x[1])
# [('H', 1), ('C', 6), ('N', 7), ('O', 8)]
\end{lstlisting}

\begin{tip}
If a lambda is complex enough that it's hard to read, use a regular function instead. Lambdas are for simple operations only.
\end{tip}

% ----------------------------------------------------------------------------
\section{Built-in Functions Worth Knowing}
\label{sec:builtin-functions}

Python includes many useful built-in functions. Here are some important ones:

\begin{lstlisting}
# Math operations on sequences
numbers = [3, 1, 4, 1, 5, 9, 2, 6]
sum(numbers)      # 31
min(numbers)      # 1
max(numbers)      # 9
len(numbers)      # 8

# Rounding
round(3.14159, 2)  # 3.14
abs(-5)            # 5

# Type conversions
int("42")          # 42
float("3.14")      # 3.14
str(42)            # "42"
list("abc")        # ['a', 'b', 'c']

# Checking types
isinstance(42, int)         # True
isinstance("hi", (int, str))  # True (either type)

# Working with sequences
sorted([3, 1, 4])           # [1, 3, 4]
reversed([1, 2, 3])         # an iterator which returns: 3, 2, 1
list(zip([1,2], ['a','b'])) # [(1,'a'), (2,'b')]
list(enumerate(['a','b']))  # [(0,'a'), (1,'b')]

# All/any
all([True, True, False])    # False (not all True)
any([True, False, False])   # True (at least one True)
\end{lstlisting}

% ----------------------------------------------------------------------------
\section{A Complete Example}
\label{sec:function-example}

Let's build a set of functions for basic spectroscopy calculations:

\begin{lstlisting}

import math # load the math library, this gives us access to logarithms by name.

def absorbance_to_transmittance(A):
    """Convert absorbance to percent transmittance."""
    return 100 * (10 ** (-A))

def transmittance_to_absorbance(T):
    """Convert percent transmittance to absorbance. This function assumes transmittence is supplied as a percentage."""
    return 2-math.log10(T)

def beers_law_concentration(A, epsilon, path_length=1.0):
    """
    Calculate concentration using Beer's Law: A = epsilon * c * l

    Parameters:
        A: Absorbance (unitless)
        epsilon: Molar absorptivity (L/(mol*cm))
        path_length: Path length in cm (default 1.0)

    Returns:
        Concentration in mol/L
    """
    return A / (epsilon * path_length)

def beers_law_absorbance(concentration, epsilon, path_length=1.0):
    """Calculate expected absorbance using Beer's Law."""
    return epsilon * concentration * path_length

# Using the functions
A_measured = 0.523
epsilon_450nm = 18500  # L/(mol*cm) for beta-carotene at 450 nm

conc = beers_law_concentration(A_measured, epsilon_450nm)
print(f"Concentration: {conc:.2e} M")

T = absorbance_to_transmittance(A_measured)
print(f"Transmittance: {T:.1f}%")
\end{lstlisting}

% ----------------------------------------------------------------------------
\section{Chapter Summary}
\label{sec:functions-summary}

Functions are fundamental to good programming:

\begin{itemize}
    \item \textbf{Define:} \py{def function_name(parameters):}
    \item \textbf{Return values:} \py{return result}
    \item \textbf{Default parameters:} \py{def func(a, b=1):}
    \item \textbf{Keyword arguments:} \py{func(a=1, b=2)}
    \item \textbf{Docstrings:} Document what functions do
    \item \textbf{Scope:} Variables inside functions are local
    \item \textbf{Lambda:} \py{lambda x: x ** 2} for simple functions
\end{itemize}

\begin{ponder}
\begin{enumerate}
    \item What's the difference between \py{print()} and \py{return}?
    \item Why should you avoid modifying global variables in functions?
    \item When would you use a lambda instead of a regular function?
    \item Write a function that takes a list of numbers and returns both the mean and standard deviation.
\end{enumerate}
\end{ponder}
